\chapter{Interface utilisateur}
\label{ch:interface}

\section{Navigation et écrans principaux}

Une barre de navigation en bas de l'écran donne accès aux cinq onglets. Avec
la fiche recette, ce sont les six écrans principaux de l'application. Les
autres écrans (mode cuisine, fiche cocktail, composition de menu, réglages)
s'ouvrent par-dessus et se ferment avec le retour système.

\begin{tableau}{L{2.6cm} Y R{1.8cm}}
  \entete{Écran} & \entete{Rôle} & \entete{Maquettes} \\
  \midrule
  Carte & Repas du jour et raccourcis de découverte & 1, 2 \\
  \midrule
  Explorer & Recherche et filtres : plats, cocktails, vins & 11, 12, 14 \\
  \midrule
  Accords & Table dressée autour d'un plat & 5 \\
  \midrule
  Courses & Liste par rayon, à cocher en magasin & 17 \\
  \midrule
  Favoris & Plats, boissons et menus enregistrés & 19 \\
  \midrule
  Fiche recette & Ingrédients, étapes, accès au mode cuisine & 3, 4 \\
\end{tableau}

\begin{planche}[fig:onglets]{Les cinq onglets de la barre de navigation}
  \maquette{1 · Accueil, la carte du jour.png}{1}{Carte du jour}
  \maquette{11 · Explorer, plats.png}{11}{Explorer}
  \maquette{5 · Accords, la table dressée.png}{5}{Accords}
  \maquette{17 · Liste de courses.png}{17}{Courses}
  \maquette{19 · Favoris.png}{19}{Favoris}
\end{planche}

\section{Parcours}

\subsection{Premier lancement}

Écran d'accueil, vérification d'âge, puis trois étapes de préférences. Chaque
étape peut être passée.

\begin{planche}[fig:lancement]{Premier lancement : âge et préférences}
  \maquette{6 · Splash.png}{6}{Splash}
  \maquette{7 · Vérification d'âge.png}{7}{Vérification d'âge}
  \maquette{8 · Onboarding, régime.png}{8}{Régime}
  \maquette{9 · Onboarding, matériel.png}{9}{Matériel}
  \maquette{10 · Onboarding, boissons.png}{10}{Boissons}
\end{planche}

\subsection{De la carte du jour à la recette}

Depuis l'accueil, l'utilisateur ouvre le plat du jour, ajuste les portions,
coche ses ingrédients puis lance le mode cuisine.

\begin{planche}[fig:recette]{Carte du jour, fiche recette et mode cuisine}
  \maquette{2 · Accueil, en défilant.png}{2}{Accueil, suite}
  \maquette{3 · Recette.png}{3}{Recette}
  \maquette{3b · Recette, Nutri-Score du plat.png}{3b}{Nutri-Score}
  \maquette{4 · Recette, ingrédients.png}{4}{Ingrédients}
  \maquette{15 · Mode cuisine.png}{15}{Mode cuisine}
\end{planche}

\subsection{Exploration}

\begin{planche}[fig:exploration]{Exploration des cocktails, des vins et des cuisines du monde}
  \maquette{12 · Explorer, cocktails.png}{12}{Cocktails}
  \maquette{13 · Cocktail, verre recommandé.png}{13}{Fiche cocktail}
  \maquette{14 · Explorer, vins.png}{14}{Vins}
  \maquette{20 · Carte du monde.png}{20}{Carte du monde}
\end{planche}

\subsection{Menu et courses}

Un menu composé se transforme en liste de courses. Le scan en magasin
(écran 18) fait partie des fonctionnalités optionnelles.

\begin{planche}[fig:menu-scan]{Composition d'un menu et scan en magasin}
  \maquette{16 · Menu builder.png}{16}{Composer le menu}
  \maquette{18 · Scan, fiche produit.png}{18}{Scan produit}
\end{planche}

\subsection{Réglages}

\begin{planche}[fig:reglages]{Réglages et sources}
  \maquette{21 · Réglages.png}{21}{Réglages}
  \maquette{22 · À propos, sources.png}{22}{À propos}
\end{planche}

\section{Composants réutilisables}

Les écrans reposent sur un petit nombre de composants communs, définis une
seule fois dans un paquet \texttt{designsystem} et partagés par toutes les
fonctionnalités.

\begin{tableau}{L{3.4cm} Y R{1.8cm}}
  \entete{Composant} & \entete{Usage} & \entete{Écrans} \\
  \midrule
  Bandeau daté & Date et météo entre un filet épais et un filet fin & 1 \\
  \midrule
  Ligne numérotée & Élément relié à sa valeur par des pointillés & 1, 5, 20, 21 \\
  \midrule
  En-tête de fiche & Origine, numéro de fiche, préparation, repos, matériel & 3 \\
  \midrule
  Étagère d'ingrédients & Vignettes posées sur un filet, cochables au toucher & 3 \\
  \midrule
  Ligne d'ingrédient & Nom, quantité et case à cocher & 4, 17 \\
  \midrule
  Boutons & Action principale pleine, secondaire en contour, compteur de portions & 3, 4, 5 \\
  \midrule
  Badges & Alcool, sans alcool, adapté, matériel & 11, 12, 13 \\
  \midrule
  Puce pays & Drapeau en tuile et nom, sélectionnable & 2, 20 \\
  \midrule
  Barre de navigation & Cinq onglets, repère sur l'onglet actif & tous \\
  \midrule
  État vide & Message et action quand une liste est vide & 19 \\
\end{tableau}

\section{Charte graphique}

Trois directions ont été explorées à partir du même écran d'accueil :
Épicerie, Cave et Marché. Nous retenons la direction C, Marché : titres en
Schibsted Grotesk, texte et chiffres en Onest, palette verte et agrume
déclinée en clair et en sombre (figure~\ref{fig:direction-c}).

Les maquettes de ce document ont été réalisées avec la première version du
design system. La direction C en garde la structure et les composants ; seuls
les couleurs, les polices et les numéros changent. Comme ces valeurs sont
centralisées dans le thème Compose, passer d'une direction à l'autre ne
modifie aucun écran.

\begin{figure}[htbp]
  \centering
  \includegraphics[width=\linewidth]{Direction C — Marché.png}
  \caption{Direction C « Marché » : typographie, palette claire et sombre}
  \label{fig:direction-c}
\end{figure}

\section{Thème, tailles d'écran et accessibilité}

\paragraph{Clair et sombre} Le thème suit celui du système par défaut et se
force dans les réglages. Chaque couleur du design system existe en deux
valeurs. Toutes les maquettes en mode sombre sont en
annexe~\ref{ann:sombre}.

\paragraph{Tailles d'écran} Les maquettes sont dessinées en 390 × 844~dp.
Les mises en page utilisent des largeurs relatives et s'adaptent aux petits
comme aux grands téléphones. Le mode cuisine reste utilisable en paysage.
Les tablettes ne sont pas visées.

\paragraph{Accessibilité} Les textes sont en sp et suivent la taille de police
du système. Les zones tactiles font au moins 48~dp. Les animations sont
réduites si le système le demande, et la photo n'est jamais la seule porteuse
d'une information.
